\documentclass[11pt]{article}
\usepackage[ngerman]{babel}
\usepackage[a4paper,margin=0.9in]{geometry}
\usepackage[T1]{fontenc}
\usepackage[utf8]{inputenc}
\usepackage{lmodern}
\usepackage{amsmath,amssymb,amsthm,mathtools,mathrsfs}
\usepackage{microtype}
\usepackage{fancyhdr}
\usepackage{array,longtable,enumitem,multicol}
\allowdisplaybreaks
\setlength{\parindent}{1.5em}
\setlength{\parskip}{0.18em}
\emergencystretch=8em
\pagestyle{fancy}
\fancyhf{}
\cfoot{\thepage}
\lhead{Heinrich Weber, Lehrbuch der Algebra}
\rhead{German source transcription}
\providecommand{\ii}{\mathrm{i}}
\providecommand{\dd}{\,d}
\begin{document}
\begin{center}{\large Erster Band. Erster Abschnitt.}\\[0.4em]{\Large \S\S~1--4. Rationale Functionen.}\end{center}

\section*{Erster Abschnitt. Rationale Functionen.}

\section*{\S~1. Ganze Functionen.}

Nächster Gegenstand der Betrachtung sind ganze rationale oder auch kurz ganze Functionen einer Veränderlichen. Wir verstehen darunter Ausdrücke von folgender Form:
\begin{equation}
 f(x)=a_0x^n+a_1x^{n-1}+a_2x^{n-2}+\cdots+a_{n-1}x+a_n,
\tag{1}
\end{equation}
worin $n$ ein ganzzahliger Exponent, der Grad der Function $f(x)$ ist. Der Grad ist eine natürliche Zahl. Bisweilen ist aber auch nützlich, von ganzen rationalen Functionen $0$ten Grades zu sprechen, worunter eine Constante verstanden wird. $x$ heisst die Veränderliche, $a_0,a_1,a_2\ldots a_{n-1},a_n$ die Coefficienten. Sowohl $x$ als $a_0,a_1,\ldots,a_n$ sind Symbole für unbestimmte Grössen, mit denen nach den Regeln der Buchstabenrechnung verfahren wird, für die auch unter Umständen bestimmte Zahlwerthe gesetzt werden können (vgl. die Einleitung).

Wenn die Function $f(x)$ in der Weise wie in (1) geschrieben ist, so nennen wir sie nach absteigenden Potenzen von $x$ geordnet. Die Summanden können in jeder beliebigen anderen Reihenfolge angeordnet, also z. B. auch nach aufsteigenden Potenzen von $x$ geordnet sein:
\[
 f(x)=a_n+a_{n-1}x+\cdots+a_1x^{n-1}+a_0x^n.
\]
Durch Addition (Subtraction) und Multiplication ganzer rationaler Functionen entstehen wieder ganze rationale Functionen. Die Vorschriften der Buchstabenrechnung geben unmittelbar die Bildungsgesetze dieser neuen Functionen.

Bei der Addition ist der Coefficient irgend einer Potenz $x^\nu$ in der Summe gleich der Summe der Coefficienten von $x^\nu$ in den einzelnen Summanden. Der Grad der entstandenen Function ist gleich dem höchsten der Grade der Summanden und kann sich nur in dem besonderen Falle erniedrigen, wenn der höchste Grad in mehreren Summanden vorkommt und die Summe der Coefficienten der höchsten Potenzen gleich Null ist.

Bei der Multiplication zweier oder mehrerer ganzer rationaler Functionen entsteht eine Function, deren Grad gleich der Summe der Grade der Factoren ist.

Um für das Product das Bildungsgesetz der Coefficienten zu übersehen, setzen wir
\begin{equation}
\begin{aligned}
 A(x)&=a_0x^m+a_1x^{m-1}+a_2x^{m-2}+\cdots,\\
 B(x)&=b_0x^n+b_1x^{n-1}+b_2x^{n-2}+\cdots,
\end{aligned}
\tag{2}
\end{equation}
\begin{equation}
 A(x)B(x)=C(x)=c_0x^{m+n}+c_1x^{m+n-1}+c_2x^{m+n-2}+\cdots
\tag{3}
\end{equation}
und erhalten
\[
 c_0=a_0b_0,
\]
\[
 c_1=a_0b_1+a_1b_0,
\]
\[
 c_2=a_0b_2+a_1b_1+a_2b_0,
\]
\[
 \cdots,
\]
oder in allgemeinen Zeichen, wenn $\nu$ eine der Zahlen $0,1,2,\ldots,m+n$ bedeutet:
\begin{equation}
 c_\nu=a_0b_\nu+a_1b_{\nu-1}+a_2b_{\nu-2}+\cdots+a_{\nu-1}b_1+a_\nu b_0,
\tag{4}
\end{equation}
worin alle $a$, deren Index grösser als $m$, und alle $b$, deren Index grösser als $n$ ist, gleich Null zu setzen sind. Denn $c_\nu$ ist der Coefficient von $x^{m+n-\nu}$, und es entsteht also $c_\nu x^{m+n-\nu}$ durch Multiplication aller Glieder von der Form $a_\mu x^{m-\mu}$ mit allen Gliedern der Form $b_{\nu-\mu}x^{n-\nu+\mu}$ und darauf folgende Addition. Sind $a_0$ und $b_0=1$, so ist auch $c_0=1$.

Hat in allen Factoren eines solchen Products die höchste Potenz von $x$ den Coefficienten $1$, so ist auch im Product der Coefficient der höchsten Potenz $=1$.

Aus diesen Vorschriften für die Rechnung mit ganzen Functionen folgt, dass die Regeln des Rechnens, die sich in Formeln wie $ab=ba$, $(ab)c=a(bc)$, $(a+b)c=ac+bc$ und ähnlichen aussprechen, auch wenn $a,b,c$ ganze Functionen von $x$ sind, gelten, und zwar in dem Sinne, dass ganze Functionen nur dann als gleich gelten, wenn gleich hohe Potenzen gleiche Coefficienten haben.

Unter ganzen Functionen mehrerer Veränderlichen verstehen wir Summen von Producten von ganzen, positiven Potenzen der Veränderlichen $x,y,z\ldots$ mit irgend welchen Coefficienten, die als constant gelten. Man kann sich diese Functionen entstanden denken aus den Functionen einer Veränderlichen $x$, wenn man darin die Coefficienten $a_0,a_1\ldots a_n$ selbst wieder als ganze rationale Functionen von anderen Veränderlichen $y,z\ldots$ auffasst. So entstehen Functionen von $m$ Veränderlichen aus Functionen von $m-1$ Veränderlichen. Alle Glieder einer solchen Function sind von der Form $x^r y^s z^t\cdots$, multiplicirt mit einem Coefficienten, und wenn die Function gehörig zusammengefasst ist, so kommt jede Combination der Exponenten $r,s,t\ldots$ nur einmal vor. Zwei so geordnete ganze Functionen gelten nur dann als einander gleich, wenn sie dieselben Producte $x^r y^s z^t\cdots$ mit denselben Coefficienten enthalten.

\section*{\S~2. Ein Satz von Gauss.}

Wir wollen sogleich eine Anwendung der Multiplicationsregel zweier ganzer rationaler Functionen machen zum Beweise eines Satzes von Gauss, der uns später noch nützlich sein wird, hier aber zur Einführung in die Rechnungsweise und als Beispiel dienen soll.\footnote{Gauss, \emph{Disquisitiones arithmeticae}, Art. 42.}

Wir betrachten hier den Fall, dass die Coefficienten in den Functionen $A(x),B(x)$ ganze Zahlen sind, so dass nach (4) auch die Coefficienten von $C(x)=A(x)B(x)$ ganze Zahlen sind.

Wenn die sämmtlichen ganzzahligen Coefficienten $a_0,a_1,\ldots,a_m$ einer Function $A(x)$ keinen gemeinschaftlichen Theiler haben, so heisst die Function $A(x)$ eine ursprüngliche oder primitive Function und der Satz, den wir beweisen wollen, lautet:

Wenn $A(x)$ und $B(x)$ ursprüngliche Functionen sind, so ist auch ihr Product $C(x)$ eine ursprüngliche Function.

Der Beweis ergiebt sich fast unmittelbar aus dem Anblick der Formel (4).

Wenn nämlich die sämmtlichen Coefficienten $c_0,c_1,c_2\ldots c_{m+n}$, wie sie in (4) angegeben sind, einen gemeinschaftlichen Theiler haben, der grösser als $1$ ist, so muss es auch wenigstens eine Primzahl geben, die in allen diesen Coefficienten aufgeht. Es sei $p$ eine solche Primzahl; diese kann nach der Voraussetzung, dass $A(x),B(x)$ ursprünglich seien, weder in allen $a_0,a_1\ldots a_m$, noch in allen $b_0,b_1\ldots b_n$ aufgehen.

Es möge nun $p$ aufgehen in
\[
 a_0,a_1\ldots a_{r-1},\quad \hbox{aber nicht in }a_r,
\]
in
\[
 b_0,b_1\ldots b_{s-1},\quad \hbox{aber nicht in }b_s,
\]
dann ist nach Voraussetzung $r\le m$ und kann auch gleich Null sein, wenn $p$ schon in $a_0$ nicht aufgeht. Ebenso ist $s\le n$.

Bilden wir nun nach (4) $c_{r+s}$ und ordnen es in folgender Weise
\begin{equation}
\begin{aligned}
 c_{r+s}={}&a_rb_s+a_{r-1}b_{s+1}+a_{r-2}b_{s+2}+\cdots\\
&+a_{r+1}b_{s-1}+a_{r+2}b_{s-2}+\cdots,
\end{aligned}
\tag{1}
\end{equation}
so sieht man unmittelbar, dass $c_{r+s}$ nicht durch $p$ theilbar sein kann, wie doch angenommen war; denn das erste Glied $a_rb_s$ ist durch $p$ nicht theilbar, während alle anderen Glieder, da sie mit einem der Coefficienten $a_0,a_1\ldots a_{r-1},b_0,b_1\ldots b_{s-1}$ multiplicirt sind, durch $p$ theilbar sind. Die Annahme also, dass $C(x)$ nicht ursprünglich sei, während es $A(x)$ und $B(x)$ sind, führt zu einem Widerspruch.

Dieser Satz lässt sich übertragen auf Functionen von mehreren Veränderlichen. Wir nennen eine ganze rationale Function von $m$ Veränderlichen mit ganzzahligen Coefficienten ursprünglich oder primitiv, wenn die Coefficienten keinen gemeinsamen Theiler haben, und wir sprechen den Satz aus, dass das Product von zwei ursprünglichen Functionen wieder eine ursprüngliche Function ist.

Um seine Wahrheit einzusehen, brauchen wir nur in der oben durchgeführten Betrachtung die Coefficienten $a_0,a_1,a_2\ldots$, $b_0,b_1,b_2\ldots$ nicht als ganze Zahlen, sondern als ganze rationale Functionen von $m-1$ Veränderlichen $y,z\ldots$ anzunehmen und eine solche Function durch eine Primzahl $p$ theilbar zu nennen, wenn alle ihre Coefficienten durch $p$ theilbar sind. Setzen wir dann voraus, der zu beweisende Satz sei bereits für Functionen von $m-1$ Variablen bewiesen, dann ergiebt die Formel (1) seine Richtigkeit für Functionen von $m$ Variablen, also seine allgemeine Gültigkeit durch den Schluss der vollständigen Induction oder den Schluss von $m-1$ auf $m$.

Eine imprimitive Function ist eine solche, deren ganzzahlige Coefficienten alle einen gemeinsamen Theiler haben. Diesen Theiler nennen wir den Theiler der Function. Dann können wir den bewiesenen Satz auch so aussprechen:

Der Theiler eines Productes zweier ganzer Functionen ist gleich dem Product der Theiler beider Functionen.

Denn sind $PA$ und $QB$ zwei ganze Functionen mit den Theilern $P$ und $Q$, so sind $A$ und $B$ ursprüngliche Functionen. Also ist auch $AB=C$ eine ursprüngliche Function und $PQ$ ist der Theiler der Function $PA\cdot QB=PQC$.

Wir können dem Satze, insofern er sich auf Functionen einer Veränderlichen bezieht, ohne seinen Inhalt wesentlich zu ändern, folgende Fassung geben, in der er besonders nützlich ist. Sind
\[
 \varphi(x)=x^m+\alpha_1x^{m-1}+\alpha_2x^{m-2}+\cdots+\alpha_m,
\]
\[
 \psi(x)=x^n+\beta_1x^{n-1}+\beta_2x^{n-2}+\cdots+\beta_n
\]
zwei ganze rationale Functionen, in denen die höchsten Potenzen von $x$ den Coefficienten $1$ haben, während die übrigen Coefficienten rationale Zahlen sind, so können in dem Product
\[
 \varphi(x)\psi(x)=x^{m+n}+\gamma_1x^{m+n-1}+\gamma_2x^{m+n-2}+\cdots+\gamma_{m+n}
\]
die Coefficienten $\gamma$ nicht alle ganze Zahlen sein, wenn die Coefficienten $\alpha,\beta$ in $\varphi(x)$ und $\psi(x)$ nicht alle ganze Zahlen sind.

Denn bezeichnen wir den kleinsten Hauptnenner der Coefficienten $\alpha$ von $\varphi$ mit $a_0$, der Coefficienten $\beta$ von $\psi$ mit $b_0$, so sind $a_0\varphi(x)=A(x)$, $b_0\psi(x)=B(x)$ primitive Functionen von $x$. Ihr Product $a_0b_0\varphi(x)\psi(x)$ hätte, wenn die Coefficienten $\gamma$ ganze Zahlen wären, den Theiler $a_0b_0$ und wäre also, wenn $a_0$ und $b_0$ nicht beide gleich $1$ wären, nicht primitiv; dies aber wäre ein Widerspruch mit dem oben bewiesenen Satze.

\section*{\S~3. Division.}

Es seien, wie bisher,
\begin{equation}
\begin{aligned}
 A&=A(x)=a_0x^m+a_1x^{m-1}+\cdots,\\
 B&=B(x)=b_0x^n+b_1x^{n-1}+\cdots
\end{aligned}
\tag{1}
\end{equation}
zwei ganze rationale Functionen von $x$; es soll aber jetzt vorausgesetzt werden, dass $m\ge n$ sei und dass $a_0$ und $b_0$ von Null verschieden sind. Dann ist die Differenz
\begin{equation}
 A-{a_0\over b_0}x^{m-n}B
\tag{2}
\end{equation}
auch eine ganze rationale Function von $x$, deren Grad aber kleiner ist als $m$, da die höchste Potenz in beiden Gliedern der Differenz denselben Coefficienten hat und also herausfällt. Wir setzen diese Differenz:
\begin{equation}
 A'=A'(x)=a'_0x^{m'}+a'_1x^{m'-1}+\cdots,
 \qquad m'<m.
\tag{3}
\end{equation}
Ist nun $m'$ noch $\ge n$, so können wir in (2) $A'$ an Stelle von $A$ setzen und dieselbe Schlussweise wiederholen.

So ergiebt sich eine Kette von Gleichungen:
\begin{equation}
\begin{aligned}
 A-{a_0\over b_0}x^{m-n}B&=A',\\
 A'-{a'_0\over b_0}x^{m'-n}B&=A'',\\
 A''-{a''_0\over b_0}x^{m''-n}B&=A''',\\
 &\cdots,
\end{aligned}
\tag{4}
\end{equation}
und diese Kette lässt sich so lange fortsetzen, bis der Grad der entstandenen Function kleiner als $n$ geworden ist. Da nun in der Reihe der Functionen $A,A',A''\ldots$ der Grad jeder folgenden mindestens um eine Einheit erniedrigt ist, so besteht die Kette der Gleichungen (4) höchstens aus $m-n+1$ Gliedern; sie kann aber auch weniger Glieder enthalten, wenn sich gleichzeitig mehrere Potenzen herausheben. Addiren wir die sämmtlichen Gleichungen (4), bezeichnen die letzte der Functionen $A,A'\ldots$ mit $C$ und setzen zur Abkürzung
\begin{equation}
 Q={a_0\over b_0}x^{m-n}+{a'_0\over b_0}x^{m'-n}+\cdots,
\tag{5}
\end{equation}
so dass auch $Q$ eine ganze rationale Function von $x$ ist, so folgt
\begin{equation}
 A=QB+C.
\tag{6}
\end{equation}
Die hier geschilderte Operation, durch die aus $A,B$ die Functionen $Q,C$ gefunden werden, heisst Division. $A$ ist der Dividendus, $B$ der Divisor, $C$ der Rest und $Q$ der Quotient. Der Grad des Restes ist immer niedriger als der Grad des Divisors.

Die Coefficienten der Functionen $Q$ und $C$ sind aus den Coefficienten $a$ und $b$ durch Addition, Subtraction, Multiplication und Theilung zusammengesetzt. Im Nenner kommen aber nur Potenzen von $b_0$ vor, und wenn also $b_0=1$ ist, so sind die Coefficienten von $Q$ und $C$ ganze Functionen der $a$ und $b$. Die höchste Potenz von $b_0$, die im Nenner auftreten kann, ist die $(m-n+1)$te, da in der Kette (4) in jeder folgenden Gleichung im Nenner einmal der Factor $b_0$ hinzukommt. Es kann aber in besonderen Fällen die höchste Potenz von $b_0$ in allen Nennern eine niedrigere sein.

Nehmen wir z. B. für $A$ eine Function dritten Grades
\begin{equation}
 f(x)=a_0x^3+a_1x^2+a_2x+a_3
\tag{7}
\end{equation}
und für $B$ die sogenannte erste Derivirte von $f(x)$, die vom zweiten Grade ist,
\begin{equation}
 f'(x)=3a_0x^2+2a_1x+a_2,
\tag{8}
\end{equation}
so erhält man:
\begin{equation}
 Q={1\over 3}x+{a_1\over 9a_0},
\tag{9}
\end{equation}
\begin{equation}
 C={6a_0a_2-2a_1^2\over 9a_0}x+{9a_0a_3-a_1a_2\over 9a_0}.
\tag{10}
\end{equation}

Wie man die in den Gleichungen (4) vorgeschriebene Rechnung zweckmässig anordnet, darf hier aus den Elementen als bekannt vorausgesetzt werden. Wir machen auf die Analogie aufmerksam, die zwischen dieser Rechnung und der Division im dekadischen Zahlsystem besteht. Eine Function $f(x)$ stellt eine dekadisch geschriebene Zahl dar, wenn die Coefficienten $a_0,a_1\ldots$ ganze Zahlen zwischen Null (einschliesslich) und $10$ (ausschliesslich) sind, und $x=10$ gesetzt wird. Lässt man in den Coefficienten auch Zahlen zu, die grösser als $10$ sind, so kann man eine Zahl auf verschiedene Arten durch $f(x)$ darstellen. Man wendet dies bei dem Divisionsverfahren an, um auf die einfachste Weise in den Resultaten gebrochene und negative Coefficienten zu vermeiden.

\section*{\S~4. Theilung durch eine lineare Function.}

Wir wollen die allgemeine Vorschrift für die Division noch auf einen besonderen Fall anwenden. Es sei der Dividend
\begin{equation}
 f(x)=a_0x^n+a_1x^{n-1}+a_2x^{n-2}+\cdots+a_n
\tag{1}
\end{equation}
beliebig, dagegen der Divisor vom ersten Grade oder, wie man auch sagt, linear. Wir wollen auch den Coefficienten der ersten Potenz von $x$ im Divisor $=1$ voraussetzen und also den Divisor in der einfachen Form $(x-\alpha)$ annehmen, worin $\alpha$ beliebig bleibt. Der Quotient $Q$ ist in diesem Falle vom $(n-1)$ten Grade und der Rest $C$ vom $0$ten Grade, d. h. von $x$ unabhängig. Setzen wir also
\begin{equation}
 f(x)=(x-\alpha)Q+C,
\tag{2}
\end{equation}
so enthält $C$ das $x$ nicht mehr, und wenn wir
\begin{equation}
 Q=q_0x^{n-1}+q_1x^{n-2}+\cdots+q_{n-2}x+q_{n-1}
\tag{3}
\end{equation}
setzen, so folgt aus (2):
\begin{equation}
\begin{aligned}
 f(x)={}&q_0x^n+q_1x^{n-1}+\cdots+q_{n-2}x^2+q_{n-1}x+C\\
&-\alpha q_0x^{n-1}-\cdots-\alpha q_{n-3}x^2-\alpha q_{n-2}x-\alpha q_{n-1},
\end{aligned}
\tag{4}
\end{equation}
und aus der Vergleichung mit (1):
\begin{equation}
\begin{aligned}
q_0&=a_0,\\
q_1-\alpha q_0&=a_1,\\
q_2-\alpha q_1&=a_2,\\
&\cdots,\\
q_{n-1}-\alpha q_{n-2}&=a_{n-1},\\
C-\alpha q_{n-1}&=a_n.
\end{aligned}
\tag{5}
\end{equation}
Daraus erhält man
\begin{equation}
\begin{aligned}
q_0&=a_0,\\
q_1&=a_0\alpha+a_1,\\
q_2&=a_0\alpha^2+a_1\alpha+a_2,\\
&\cdots,\\
q_{n-1}&=a_0\alpha^{n-1}+a_1\alpha^{n-2}+\cdots+a_{n-1},\\
C&=a_0\alpha^n+a_1\alpha^{n-1}+\cdots+a_{n-1}\alpha+a_n=f(\alpha).
\end{aligned}
\tag{6}
\end{equation}
$C$ entsteht aus $f(x)$, wenn man $x=\alpha$ setzt, und kann also auch mit $f(\alpha)$ bezeichnet werden.

Demnach haben wir auch die Formel
\begin{equation}
 {f(x)-f(\alpha)\over x-\alpha}=Q(x),
\tag{7}
\end{equation}
worin $Q(x)$ eine ganze Function vom Grade $n-1$ ist.

Wenn wir in den Ausdrücken (6) an Stelle der unbestimmten Grösse $\alpha$ das Zeichen $x$ setzen, so entsteht daraus eine Reihe von ganzen rationalen Functionen von $x$, die wir, wenn wir der Einfachheit halber $a_0=1$ setzen, so schreiben:
\begin{equation}
\begin{aligned}
f_0&=1,\\
f_1&=x+a_1,\\
f_2&=x^2+a_1x+a_2,\\
&\cdots,\\
f_{n-1}&=x^{n-1}+a_1x^{n-2}+a_2x^{n-3}+\cdots+a_{n-1}.
\end{aligned}
\tag{8}
\end{equation}
Diese Functionen $f_0,f_1\ldots f_{n-1}$ werden uns später noch gute Dienste leisten. Für jetzt fügen wir noch folgende Bemerkungen bei.

Man kann nach (8) die Potenzen $1,x,x^2\ldots x^{n-1}$ von $x$ linear ausdrücken durch die Functionen $f_0,f_1\ldots f_{n-1}$ und zwar so, dass in den Coefficienten nur ganze rationale Verbindungen der $a$ vorkommen, z. B.
\[
 1=f_0,
\]
\[
 x=f_1-a_1f_0,
\]
\[
 x^2=f_2-a_1f_1+(a_1^2-a_2)f_0,
\]
\[
 \cdots,
\]
und daraus folgt, dass man jede ganze rationale Function von $x$, deren Grad nicht grösser als $n-1$ ist, gleichfalls linear durch $f_0,f_1,\ldots,f_{n-1}$ ausdrücken kann in der Form
\begin{equation}
 y_0f_0+y_1f_1+\cdots+y_{n-1}f_{n-1},
\tag{9}
\end{equation}
worin die Coefficienten $y_0,y_1\ldots y_{n-1}$ von $x$ unabhängig sind.

Ist also $F(x)$ eine beliebige ganze rationale Function von $x$, so kann man nach \S~3, indem man $f(x)$ als Divisor betrachtet,
\begin{equation}
 F(x)=Qf(x)+y_0f_0+y_1f_1+\cdots+y_{n-1}f_{n-1}
\tag{10}
\end{equation}
setzen, worin auch $Q$ eine ganze rationale Function von $x$ ist.

Zur recurrenten Berechnung der Functionen $f_r(x)$ ergiebt sich aus (8) die Relation:
\begin{equation}
 f_r(x)-xf_{r-1}(x)=a_r.
\tag{11}
\end{equation}

\end{document}
